\section{The staleness dual}
\label{sec:duality}

\spec{9.2} walks \emph{down} from the packages a commit changed. An audit
command wants the opposite walk: given a package, which of its dependencies is
it behind, and by how much? \spec{13.7b} gives that walk and asserts that the
two agree, calling disagreement an implementation bug. This section proves
that the audit is sound against \spec{9.2} unconditionally, that it is
complete under two conditions the current text does not state, and exhibits
the disagreements outside them. The pseudo-code is reproduced from
\spec{13.7b} unchanged, except that the arrow and ellipsis glyphs in the
comment on \code{dist}, which the monospace font lacks, are set as
\code{->} and \code{...}.

\begin{speccode}[fontsize=\footnotesize]
staleSources(D):
    edges = graph.restrictedTo(config.propagation.kinds)  # §8.4: the same edges as §9.2
    channel = resolveChannelsForAudit()                    # same phase-2 result as §9.2
    dist  = shortestPathsUp(D, edges)   # P -> edge count of the shortest path P -> ... -> D
    out   = {}
    for u in units:                                       # every retained tuple, §13.4
        if bumpOf(u) == none:                             continue
        b = (u.propagate == 'inherit') ? bumpOf(u) : u.propagate    # as §9.2
        if b == none or u.depth == 0:                     continue
        sources = sourcePackages(u)                       # §13.4a
        if sources is empty:                              continue
        if D in sources:                                  continue   # §9.2 seeds seen = sources
        sources = owed(u, D)                              # §13.4a: what D is still owed
        if sources is empty:                              continue   # every source delivered
        reaching = { P in sources : P in dist }
        if reaching is empty:                             continue
        if not resolvableBy(sources, D, channel):     continue   # §9.3a
        level = min({ dist[P] for P in reaching })        # measured from the whole source set
        if level > u.depth:                               continue
        if D not in resolve(u.propagateScope):            continue
        if cancelledFor(commitOf(u), D):                  continue
        for P in reaching:
            out |= { (P, u, level, b) }
    return out
\end{speccode}

Three things changed from the CCME 2.0.0 audit. The loop runs over every
retained tuple, not over the units of $\Winf{}{D}$: a commit $D$ has released
past can still be owed to it by a source that released it later, and iterating
the target's window alone misses exactly the consumers that got ahead of a
provider, which are the ones the audit exists to find after a partial failure.
\code{sources} is rebound to $\owed{}{u}{D}$ before the distance and
resolvability tests, so the rows name the owed sources, as \code{prov[D]} does
in \spec{9.2}. And \code{resolvableBy} is tested before the level test.

\begin{definition}
\label{def:stale}
$\stale{\st}{d} = \{(p,u,\ell,b)\}$ is the set the procedure returns. For the
rows of a unit $u$, $p$ ranges over the owed sources of $u$ that reach $d$,
$\ell = \dist{\owed{\st}{u}{d}}{d}$, the least distance from an owed source,
and $b = \pbump{u}$. Write $\srcch{\st}{u}' = \{\chanres{\st}{P} : P \in
\owed{\st}{u}{d}\}$ for the channel set the audit tests, and
$\resolv{\st}{u}{d}'$ for \cref{def:resolvable} evaluated over it.
\end{definition}

\begin{theorem}[duality]
\label{thm:duality}
For every state \st{} and package $d$:
\begin{enumerate}[label=(\alph*),leftmargin=2.4em]
  \item (soundness) if $u$ contributes a row to $\stale{\st}{d}$ then
    $\admit{\st}{u}{d}$; hence $\stale{\st}{d} \neq \emptyset$ implies
    $\propb{\st}{d} \neq \bnone$ and $\bigsqcup\{\, b : (p,u,\ell,b) \in
    \stale{\st}{d} \,\} \leq \propb{\st}{d}$; \label{du:sound}
  \item (completeness) if (R) holds and $\resolv{\st}{u}{d}' =
    \resolv{\st}{u}{d}$ for every $u$ with $\owed{\st}{u}{d} \neq \emptyset$,
    then $u$ contributes a row exactly when $\admit{\st}{u}{d}$, and
    \[
      \stale{\st}{d} \neq \emptyset
      \iff \propb{\st}{d} \neq \bnone,
      \qquad
      \bigsqcup\{\, b : (p,u,\ell,b) \in \stale{\st}{d} \,\} = \propb{\st}{d}.
    \] \label{du:complete}
\end{enumerate}
\end{theorem}

\begin{proof}
Every row of $u$ carries $b = \pbump{u}$, and $\propb{\st}{d}$ is the join of
$\pbump{u}$ over admitted units (\cref{def:propagated}), so both parts reduce
to comparing ``$u$ contributes a row'' with $\admit{\st}{u}{d}$.

\ref{du:sound}. A row means every guard passed. \code{bumpOf(u) == none},
\code{b == none} and \code{u.depth == 0} are \cref{a:live}. \code{sources is
empty} on the full source set is the first half of \cref{a:reach}, and
\code{D in sources} is the $d \notin \srcp{\st}{u}$ of \cref{def:reach}, read
from the other end: it is the same clause as seeding \code{seen} with the
sources. \code{sources is empty} after the rebinding is \cref{a:window}.
\code{level} is $\min\{\dist{\{P\}}{d} : P \in \code{reaching}\} =
\dist{\owed{\st}{u}{d}}{d}$, because \code{shortestPathsUp} gives
$\dist{\{P\}}{d}$ for each $P$ and a non-reaching source has distance
$\infty$; and $\owed{\st}{u}{d} \subseteq \srcp{\st}{u}$ gives
$\dist{\srcp{\st}{u}}{d} \leq \dist{\owed{\st}{u}{d}}{d} = \code{level} \leq
\pdepth{u}$, which with $d \notin \srcp{\st}{u}$ is the second half of
\cref{a:reach} by \cref{lem:reach}. \code{D not in resolve(u.propagateScope)}
is \cref{a:scope}; \code{cancelledFor} is \cref{a:cancel}. Finally
$\srcch{\st}{u}' \subseteq \srcch{\st}{u}$, so $\resolv{\st}{u}{d}'$ implies
$\resolv{\st}{u}{d}$, which is \cref{a:resolvable}. So every guarded unit is
admitted, and the join over a subset of the admitted units is at most the join
over all of them.

\ref{du:complete}. Conversely let $\admit{\st}{u}{d}$. \Cref{a:live},
\cref{a:reach}'s first half, \cref{a:window}, \cref{a:scope} and
\cref{a:cancel} pass the corresponding guards as above, and $d \notin
\srcp{\st}{u}$ passes \code{D in sources}. By (R) every owed source reaches
$d$, so \code{reaching} $= \owed{\st}{u}{d} \neq \emptyset$, and
$\code{level} = \dist{\owed{\st}{u}{d}}{d} = \dist{\srcp{\st}{u}}{d} \leq
\pdepth{u}$ by \cref{a:reach}. By the second hypothesis \cref{a:resolvable}
passes \code{resolvableBy} over the owed sources. So $u$ contributes a row,
the guarded set is the admitted set, non-emptiness is equivalent to
$\propb{\st}{d} \neq \bnone$, and the joins agree.
\end{proof}

\begin{remark}[where the two formulations disagree]
\label{rem:owed-reach}
\spec{13.7b} says the two formulations MUST agree and that disagreement is an
implementation bug. As written they do not agree, and the disagreement is in
the text, not in an implementation; three shapes, all on \cref{fig:workspace},
all with a multi-source unit and a target that got ahead of one source.
\begin{itemize}
  \item \emph{An owed source that does not reach the target.} Unit
    \code{feat(core,api)\^{}: x} at $C$; \code{cli} depends on \code{core}
    only. Let \code{core} have delivered $C$ to \code{cli} and \code{api} have no
    release carrying $C$. \spec{9.2} reaches \code{cli} at level $1$ through
    \code{core}, finds $\owed{}{u}{\code{cli}} = \{\code{api}\}$, and bumps
    \code{cli} with provenance \code{api}, a package it does not depend on.
    The audit rebinds \code{sources} to $\{\code{api}\}$, finds
    \code{reaching} empty, and reports nothing. \Cref{def:wf} then also demands
    that \code{api} precede \code{cli} in the publish order whenever
    \code{cli} is tagged at \Head{}, a constraint no edge supports.
  \item \emph{An owed source further away than a delivered one.} Unit
    \code{feat(core,ui)\^{}: x} at $C$, target \code{@acme/theme}, which
    depends on \code{ui} and reaches \code{core} at distance $2$. Let
    \code{ui} have delivered $C$ to \code{@acme/theme} and \code{core} not.
    \spec{9.2} admits at level $1$ with $\owed{}{u}{\code{@acme/theme}} =
    \{\code{core}\}$; the audit computes \code{level} $= 2 > 1$ and reports
    nothing.
  \item \emph{Resolvability over different sets.} The same unit with
    \code{core} on \stable{} and \code{ui} on \code{beta}, target
    \code{@acme/theme} on \stable{}, \code{core} delivered and \code{ui} owed.
    \spec{9.2} hoists \code{anyStable} over the full source set and admits;
    the audit tests $\{\code{beta}\}$ against a stable target and rejects,
    which is what \spec{9.3a} intends.
\end{itemize}
In each case the audit under-reports relative to \spec{9.2}; by
\cref{thm:duality}\ref{du:sound} it never over-reports. The first shape is the
more serious, since \spec{9.2}'s answer is the questionable one: a bump owed
by a package the target cannot reach, and a well-formedness constraint over a
non-edge. A definition of $\owed{}{u}{D}$ restricted to sources that reach $D$
would remove the first two shapes and make (R) a theorem; the third needs
\spec{9.2} and \spec{13.7b} to agree on which source set resolvability is
evaluated over. All three are flagged for the specification.

The second shape is not hypothetical: the reference engine produced it until product \code{main} \code{e461d7d9}, which came after the build the campaigns of \cref{sec:evaluation} measured; a consumer that \code{ui} had delivered a shared commit to was released again on \code{core}'s account, two edges away (\cref{ob:attribution}). The specification's text permits it; the fix narrows the engine to the reaching set.
\end{remark}

\begin{example}[the audit, on both failed runs]
\label{ex:stale}
At $\st_1$ of the first shape of \cref{ex:orphan}, where \code{core} has
published and \code{cli} has not, the audit for $d = \code{cli}$ walks up from
\code{cli}, finds \code{core} at $\code{dist}[\code{core}] = 1$, finds
$\owed{\st_1}{u}{\code{cli}} = \{\code{core}\}$, and returns the single row
\[
  \stale{\st_1}{\code{cli}} =
  \bigl\{\,(\code{core},\ u,\ 1,\ \bpatch)\,\bigr\},
\]
whose join is \bpatch{} $= \propb{\st_1}{\code{cli}}$, as \cref{thm:duality}
requires; the unit has one source, so (R) and the resolvability condition hold
outright. Run the same audit for \code{ui}, which did publish, and it is empty:
$\deliv{\st_1}{\code{core}}{C}{\code{ui}}$ empties the owed set. At $\st_1$ of
the second shape, where \code{cli} proceeded and is tagged at $C$, the audit
returns the same row for \code{cli}, although $C \notin
\Winf{\st_1}{\code{cli}}$: this is the row the CCME 2.0.0 audit, which
iterated the target's fresh window, could not produce, and it is the question
an operator actually asks after a failed run (which of my packages are behind,
and behind what), answered without recomputing a plan.
\end{example}

The clauses the proof leans on are the ones \spec{13.7b} warns an
implementation drifts on, and it is worth naming what each drift produces.
Excluding $d$ from its own unit's sources: without it, a package that both
changes and consumes a sibling in one unit audits as behind itself. Measuring
\code{level} from the whole set rather than per ancestor: taking
$\code{dist}[p]$ for each $p$ separately over-reports the level and so
under-reports staleness, exactly in the diamond cases. Iterating the target's
window rather than every retained tuple: the consumer that got ahead is
invisible. And resolving \code{inherit} to the unit's own bump before testing
it against \bnone.

\begin{proposition}[the tag-level screen is necessary, not sufficient]
\label{prop:screen}
For packages $P, D$ with baselines, say the screen of \spec{13.7b}
\emph{fires} when $\basec{\st}{P} \not\anc \basec{\st}{D}$. Then:
\begin{enumerate}[label=(\alph*),leftmargin=2.4em]
  \item the screen fires iff $\lnot\deliv{\st}{P}{\basec{\st}{P}}{D}$: it is
    \cref{def:delivery} with the unit's commit replaced by $P$'s own baseline
    commit; \label{sc:shape}
  \item if $\admit{\st}{u}{d}$, $P \in \owed{\st}{u}{d}$ and $\cof{u} \anc
    \basec{\st}{P}$ (the source has released the work and its baseline reaches
    that release), then the screen fires for $(P,d)$; \label{sc:nec}
  \item the converse of \ref{sc:nec} fails. \label{sc:conv}
\end{enumerate}
\end{proposition}

\begin{proof}
\ref{sc:shape}: $\deliv{\st}{P}{\basec{\st}{P}}{D}$ asks for $t \in
\relc{\st}{P}$ with $\basec{\st}{P} \anc t \anc \basec{\st}{D}$; any such $t$
gives $\basec{\st}{P} \anc \basec{\st}{D}$ by transitivity, and conversely $t =
\basec{\st}{P}$ is a release commit of $P$.
\ref{sc:nec}: if $\basec{\st}{P} \anc \basec{\st}{d}$ then $t = \basec{\st}{P}$
witnesses $\deliv{\st}{P}{\cof{u}}{d}$, contradicting $P \in
\owed{\st}{u}{d}$.
\ref{sc:conv}: a provider may be legitimately ahead of a consumer with
nothing owed: the units between them may all declare \code{\^{}none} or depth
\code{+0}, may be excluded by \code{Propagate-Scope}, may reach $d$ only beyond
their declared depth, or may travel edges of a kind not in
\code{propagation.kinds}. Each gives a state where the screen fires and no bump
is admitted.
\end{proof}

The screen is therefore a way to \emph{find} candidates cheaply, by comparing
two tags rather than walking a window, and never a way to decide releases;
using it to decide would manufacture bumps that no commit asked for. Part
\ref{sc:shape} is why it screens: a consumer owed something has a provider
release its own release does not reach. Part \ref{sc:nec} needs the provider
to have released the work; a debt from a provider that has not released the
commit at all is not the screen's business, since that provider is then in
the plan on its own account (\cref{lem:tagged-target}). There is also a case
the screen cannot see at all: $P$ and $d$ released at the same commit, in the
same run, with $d$ published first. Ancestry cannot distinguish them and no
rule over tags can. That case is prevented rather than detected, by publishing
in dependency order (\cref{def:order}) and by refusing a run that would
release a provider at the baseline commit of a consumer it still owes without
that consumer following (\cref{def:wf}); where the consumer follows and fails,
\cref{prop:loss} is the reported remainder.

\begin{remark}[historical source of the eligibility guard]
\label{rem:duality-hyp}
CCME 1.0.0 omitted the resolvability guard, so wherever \spec{9.3a} suppressed
a bump the two formulations disagreed. \Cref{ob:duality} preserves the
counterexample; CCME 2.0.0 made the repair normative, and the current text
keeps it while adding the delivery rebinding whose consequences
\cref{rem:owed-reach} records.
\end{remark}
